
{% Q108868

\needspace{31\baselineskip} 
\item \rtask \ponto{\pt} % 
A programação de computadores corresponde à escrita, teste e manutenção de programas de computador. Sobre a programação de computadores é CORRETO afirmar:

\begin{enumerate}[label={\Roman*.},leftmargin=*]
    \item Um algoritmo corresponde a uma sequência ordenada, e sem ambiguidade, de ações que levam à solução de um problema e, quando codificado em uma linguagem de programação, corresponde a um programa de computador.
    
    \item Recursividade é uma técnica de programação eficaz para resolver um problema originalmente complexo, reduzindo-o em pequenas ocorrências do problema principal. Um algoritmo é dito recursivo quando chama a si mesmo ou chama uma sequência de outros algoritmos, e um deles chama novamente o primeiro algoritmo.
    
    \item Um \textit{array} é uma estrutura de dados heterogênea onde seus elementos individuais são acessados através de índices que indicam sua posição na estrutura.
    
    \item Procedimentos e funções são sub-algoritmos codificados como parte do desenvolvimento de um algoritmo para a solução de um problema particular.
\end{enumerate}

A sequência correta é:

\begin{answerlist}[label={\texttt{\Alph*}.},leftmargin=*]
    \ti Apenas as assertivas II, III e IV estão corretas.
    \ti Apenas as assertivas I e IV estão corretas.
    \di Apenas as assertivas I, II e IV estão corretas.
    \ti As assertivas I, II, III e IV estão corretas.
    \ti Apenas as assertivas III e IV estão corretas.
\end{answerlist}

    % Alternativa correta: C - Apenas as assertivas I, II e IV estão corretas.
    % 
    % Vamos discorrer sobre cada assertiva para entender o porquê da alternativa C ser a correta:
    % 
    % I. Um algoritmo corresponde a uma sequência ordenada, e sem ambiguidade, de ações que levam à solução de um problema e, quando codificado em uma linguagem de programação, corresponde a um programa de computador. Esta assertiva está correta. O algoritmo é de fato uma descrição passo a passo de como resolver um determinado problema e, quando transformado em código usando uma linguagem de programação, torna-se um programa executável em um computador.
    % 
    % II. Recursividade é uma técnica de programação eficaz para resolver um problema originalmente complexo, reduzindo-o em pequenas ocorrências do problema principal. Um algoritmo é dito recursivo quando chama a si mesmo ou chama uma sequência de outros algoritmos, e um deles chama novamente o primeiro algoritmo. Esta assertiva também está correta. A recursividade é uma abordagem onde a função chama a si mesma para resolver problemas menores do mesmo tipo, o que pode muitas vezes simplificar a solução de problemas complexos.
    % 
    % III. Um array é uma estrutura de dados heterogênea onde seus elementos individuais são acessados através de índices que indicam sua posição na estrutura. Esta assertiva está incorreta. Na verdade, arrays são estruturas de dados homogêneas, o que significa que todos os seus elementos devem ser do mesmo tipo. A heterogeneidade não é característica de um array, mas sim de outras estruturas como, por exemplo, structs em C ou objetos em linguagens orientadas a objeto.
    % 
    % IV. Procedimentos e funções são sub-algoritmos codificados como parte do desenvolvimento de um algoritmo para a solução de um problema particular. Esta assertiva está correta. Procedimentos e funções são blocos de código que realizam tarefas específicas e podem ser chamados repetidamente. Eles ajudam a reduzir a repetição de código e tornam o algoritmo principal mais fácil de entender e manter.
    % 
    % Portanto, as assertivas I, II e IV estão corretas, ressaltando a importância de entender conceitos como algoritmos, recursividade, estruturas de dados e modularização na programação, enquanto a assertiva III continha um erro ao classificar um array como heterogêneo.
}






